% Capitulo 6: bibliografia con seis citas distintas.
\section{Bibliografía}\label{sec:bib}

Las citas viven en \texttt{bib/references.bib} y se resuelven al componer,
con el estilo \texttt{plain} al final del documento. Este capítulo cita
seis entradas distintas para demostrar que el mecanismo funciona desde
cualquier capítulo: tres herramientas reales que sostienen el flujo y tres
entradas históricas del ejemplo original.

El motor que lo compone todo es tectonic~\cite{tectonic}, el corazón
embebido que presentó el capítulo~\ref{sec:que-es}. Los diagramas de
secuencia hablan Mermaid~\cite{mermaid} y los grafos dirigidos hablan el
lenguaje DOT de Graphviz~\cite{graphviz}; ambos se citaron ya en el
capítulo~\ref{sec:diagramas}, donde se componen.

Junto a ellas conviven las tres entradas originales del ejemplo: un
artículo de revista~\cite{example}, un libro~\cite{example2} y un recurso
en línea~\cite{example3}. Se conservan tal cual para que los documentos
que ya las citaban sigan compilando, y sirven aquí de testigos de que las
citas clásicas --- artículo, libro y web --- siguen resolviéndose como
siempre. \texttt{texforge check} avisa de toda clave \texttt{\textbackslash
  cite} sin entrada en el \texttt{.bib} antes de compilar.

El flujo bibliográfico no tiene misterio: cada obra vive una vez en
\texttt{bib/references.bib} con su clave, se cita en el texto con
\texttt{\textbackslash cite} y aparece formateada al final, ordenada según
el estilo elegido, aquí \texttt{plain}. Las claves conviene que sean
cortas y memorables --- el apellido, la herramienta, el año --- porque se
escriben a mano en cada cita. Y cuando el documento crece, el mismo
\texttt{.bib} sirve a varios capítulos sin duplicar nada: este capítulo
cita seis entradas distintas repartidas por todo el documento, y la lista
final las reúne a todas sin que nadie las ordene a mano.

\begin{notebox}
  Las tres primeras entradas vienen del ejemplo original y se conservan a
  propósito: los documentos que ya las citaban siguen compilando sin
  cambios, y aquí demuestran que lo clásico convive con lo nuevo.
\end{notebox}
